% ============================================================================
%  第 2 章 相关工作 —— body/02-related-work.tex
%  职责：把本文放回到坐标系里，突出「别人没做的那格」。
%  写法：分类评述 → 对比表 → 明确空白
% ============================================================================

\section{相关工作}
\label{sec:related-work}

% ---------------------------------------------------------------------------
\subsection{第三方风险建模}

坠机概率模型经历了从常数假设向动态建模的演进。早期研究以恒定失效率
或经验图谱刻画 $\pcrash$ \citep{primatesta2022review}，
难以反映环境状态对失效概率的影响。近期工作引入比例风险模型处理
微气象协变量 \citep{cooper2023proportional}，但在暴露度与时间维度上
仍沿用静态假设。

% ---------------------------------------------------------------------------
\subsection{风险感知路径规划}

基于风险场的路径搜索方法可归纳为两类：栅格/概率地图上的改进启发式搜索
\citep{zhang2024spatiotemporal}，以及以可达安全性为约束的优化/采样方法
\citep{chen2024costastar}。前者多采用静态风险场，后者则通常假设
滑动窗内的短期最优。

% ---------------------------------------------------------------------------
\subsection{时间与时空依赖的搜索方法}

当边代价随通过时刻变化时，问题退化为时间依赖最短路：代价函数须满足
FIFO（先到先走）性质，标签更改/标签更正算法构成主流求解框架。
\todo{补 citation：time-dependent shortest path 经典文献。}

% ---------------------------------------------------------------------------
\subsection{研究空白定位}

\reftab{tab:literature-comparison} 从「风险维度、时间动态、暴露度演化、求解方式」
四个维度对比代表性工作，据此定位本文所处的空白。

% 表格统一放在 assets/tables/ 并以 \input 引入，便于多人协作时分离文件冲突
% ============================================================================
%  文献对比表 —— assets/tables/literature-comparison.tex
%  位置：第 \ref{sec:related-work} 节
% ============================================================================

\begin{table}[t]
  \centering
  \caption{相关工作在四个维度上的对比（$\checkmark$~支持，$	imes$~不支持，$\sim$~部分支持）}
  \label{tab:literature-comparison}
  \footnotesize
  \begin{threeparttable}
  \begin{tabular}{@{}lcccc@{}}
    \toprule
    工作 & 多维风险 & 时间动态 & 暴露度演化 & 最优性保证 \\
    \midrule
    Primatesta 等 \cite{primatesta2022review} & $\checkmark$        & $	imes$           & $	imes$           & --          \\
    Cooper 等 \cite{cooper2023proportional}   & $	imes$            & $\checkmark$       & $	imes$           & --          \\
    Zhang 等 \cite{zhang2024spatiotemporal}   & $\sim$              & $\sim$             & $	imes$           & $	imes$    \\
    Chen 等 \cite{chen2024costastar}          & $	imes$            & $\sim$             & $	imes$           & $\sim$      \\
    \midrule
    \textbf{本文} & $\checkmark$ & $\checkmark$ & $\checkmark$ & $\checkmark$ \\
    \bottomrule
  \end{tabular}
  \begin{tablenotes}
    \footnotesize
    \item[] 注：$\checkmark$/$	imes$/$\sim$ 分别表示支持、不支持、部分支持。
  \end{tablenotes}
  \end{threeparttable}
\end{table}

